\section{Replication on the 2025 Survey of Consumer Finances}
\label{app:scf2025}

% DRAFT for review. Every number below is a \rep macro generated by
% `imp replication` from the committed run directories (paper/tables/
% scf2025_numbers.tex); tests/test_replication.py fails if a number is typed
% in by hand. The rules are docs/scf-2025-replication-criteria.md, committed
% before any method was fit on the 2025 wave.
% Generated by `imp replication` -- do not edit by hand. Every number the SCF 2025 appendix cites, as \rep{key}.
\providecommand{\rep}[1]{\ifcsname rep:#1\endcsname\csname rep:#1\endcsname\else\textbf{[missing: #1]}\fi}
\expandafter\def\csname rep:F0:status2022\endcsname{met}
\expandafter\def\csname rep:F0:status2025\endcsname{met}
\expandafter\def\csname rep:F0:verdict\endcsname{replicated}
\expandafter\def\csname rep:F1a:parts\endcsname{6}
\expandafter\def\csname rep:F1a:parts-met:2022\endcsname{6}
\expandafter\def\csname rep:F1a:parts-met:2025\endcsname{5}
\expandafter\def\csname rep:F1a:status2022\endcsname{met}
\expandafter\def\csname rep:F1a:status2025\endcsname{inconclusive}
\expandafter\def\csname rep:F1a:verdict\endcsname{inconclusive}
\expandafter\def\csname rep:F1b:status2022\endcsname{met}
\expandafter\def\csname rep:F1b:status2025\endcsname{met}
\expandafter\def\csname rep:F1b:verdict\endcsname{replicated}
\expandafter\def\csname rep:F1c:status2022\endcsname{met}
\expandafter\def\csname rep:F1c:status2025\endcsname{met}
\expandafter\def\csname rep:F1c:verdict\endcsname{replicated in direction, magnitude differs (smaller)}
\expandafter\def\csname rep:F2a:parts\endcsname{1}
\expandafter\def\csname rep:F2a:parts-met:2022\endcsname{1}
\expandafter\def\csname rep:F2a:parts-met:2025\endcsname{1}
\expandafter\def\csname rep:F2a:status2022\endcsname{met}
\expandafter\def\csname rep:F2a:status2025\endcsname{met}
\expandafter\def\csname rep:F2a:verdict\endcsname{replicated}
\expandafter\def\csname rep:F2b:status2022\endcsname{met}
\expandafter\def\csname rep:F2b:status2025\endcsname{met}
\expandafter\def\csname rep:F2b:verdict\endcsname{replicated}
\expandafter\def\csname rep:F2c:parts\endcsname{4}
\expandafter\def\csname rep:F2c:parts-met:2022\endcsname{4}
\expandafter\def\csname rep:F2c:parts-met:2025\endcsname{3}
\expandafter\def\csname rep:F2c:status2022\endcsname{met}
\expandafter\def\csname rep:F2c:status2025\endcsname{failed}
\expandafter\def\csname rep:F2c:verdict\endcsname{not replicated}
\expandafter\def\csname rep:F2d:status2022\endcsname{met}
\expandafter\def\csname rep:F2d:status2025\endcsname{met}
\expandafter\def\csname rep:F2d:verdict\endcsname{replicated}
\expandafter\def\csname rep:F3a:status2022\endcsname{met}
\expandafter\def\csname rep:F3a:status2025\endcsname{met}
\expandafter\def\csname rep:F3a:verdict\endcsname{replicated}
\expandafter\def\csname rep:F3b:status2022\endcsname{met}
\expandafter\def\csname rep:F3b:status2025\endcsname{met}
\expandafter\def\csname rep:F3b:verdict\endcsname{replicated}
\expandafter\def\csname rep:F3c:status2022\endcsname{met}
\expandafter\def\csname rep:F3c:status2025\endcsname{failed}
\expandafter\def\csname rep:F3c:verdict\endcsname{not replicated}
\expandafter\def\csname rep:F4a:parts\endcsname{4}
\expandafter\def\csname rep:F4a:parts-met:2022\endcsname{4}
\expandafter\def\csname rep:F4a:parts-met:2025\endcsname{3}
\expandafter\def\csname rep:F4a:status2022\endcsname{met}
\expandafter\def\csname rep:F4a:status2025\endcsname{inconclusive}
\expandafter\def\csname rep:F4a:verdict\endcsname{inconclusive}
\expandafter\def\csname rep:F4b:parts\endcsname{2}
\expandafter\def\csname rep:F4b:parts-met:2022\endcsname{2}
\expandafter\def\csname rep:F4b:parts-met:2025\endcsname{2}
\expandafter\def\csname rep:F4b:status2022\endcsname{met}
\expandafter\def\csname rep:F4b:status2025\endcsname{met}
\expandafter\def\csname rep:F4b:verdict\endcsname{replicated}
\expandafter\def\csname rep:F4c:status2022\endcsname{met}
\expandafter\def\csname rep:F4c:status2025\endcsname{failed}
\expandafter\def\csname rep:F4c:verdict\endcsname{not replicated}
\expandafter\def\csname rep:F5a:status2022\endcsname{met}
\expandafter\def\csname rep:F5a:status2025\endcsname{met}
\expandafter\def\csname rep:F5a:verdict\endcsname{replicated}
\expandafter\def\csname rep:F5b:parts\endcsname{4}
\expandafter\def\csname rep:F5b:parts-met:2022\endcsname{4}
\expandafter\def\csname rep:F5b:parts-met:2025\endcsname{2}
\expandafter\def\csname rep:F5b:status2022\endcsname{met}
\expandafter\def\csname rep:F5b:status2025\endcsname{failed}
\expandafter\def\csname rep:F5b:verdict\endcsname{not replicated}
\expandafter\def\csname rep:G1:status2022\endcsname{met}
\expandafter\def\csname rep:G1:status2025\endcsname{failed}
\expandafter\def\csname rep:G1:verdict\endcsname{not replicated}
\expandafter\def\csname rep:W1:status2022\endcsname{met}
\expandafter\def\csname rep:W1:status2025\endcsname{met}
\expandafter\def\csname rep:W1:verdict\endcsname{replicated}
\expandafter\def\csname rep:W2:status2022\endcsname{met}
\expandafter\def\csname rep:W2:status2025\endcsname{met}
\expandafter\def\csname rep:W2:verdict\endcsname{replicated in direction, magnitude differs (smaller)}
\expandafter\def\csname rep:W3a:status2022\endcsname{met}
\expandafter\def\csname rep:W3a:status2025\endcsname{met}
\expandafter\def\csname rep:W3a:verdict\endcsname{replicated}
\expandafter\def\csname rep:W3b:status2022\endcsname{inconclusive}
\expandafter\def\csname rep:W3b:status2025\endcsname{inconclusive}
\expandafter\def\csname rep:W3b:verdict\endcsname{inconclusive}
\expandafter\def\csname rep:W3c:status2022\endcsname{inconclusive}
\expandafter\def\csname rep:W3c:status2025\endcsname{met}
\expandafter\def\csname rep:W3c:verdict\endcsname{inconclusive}
\expandafter\def\csname rep:W4a:status2022\endcsname{met}
\expandafter\def\csname rep:W4a:status2025\endcsname{met}
\expandafter\def\csname rep:W4a:verdict\endcsname{replicated}
\expandafter\def\csname rep:W4b:status2022\endcsname{met}
\expandafter\def\csname rep:W4b:status2025\endcsname{met}
\expandafter\def\csname rep:W4b:verdict\endcsname{replicated}
\expandafter\def\csname rep:W5:status2022\endcsname{met}
\expandafter\def\csname rep:W5:status2025\endcsname{met}
\expandafter\def\csname rep:W5:verdict\endcsname{replicated}
\expandafter\def\csname rep:W6:status2022\endcsname{met}
\expandafter\def\csname rep:W6:status2025\endcsname{met}
\expandafter\def\csname rep:W6:verdict\endcsname{replicated}
\expandafter\def\csname rep:W7:status2022\endcsname{met}
\expandafter\def\csname rep:W7:status2025\endcsname{met}
\expandafter\def\csname rep:W7:verdict\endcsname{replicated}
\expandafter\def\csname rep:check:families:2022\endcsname{4,595}
\expandafter\def\csname rep:check:families:2025\endcsname{4,360}
\expandafter\def\csname rep:check:net_worth_mean:gap:2022\endcsname{$-$0.40}
\expandafter\def\csname rep:check:net_worth_mean:gap:2025\endcsname{$-$0.77}
\expandafter\def\csname rep:check:net_worth_mean:public:2022\endcsname{1,159.4}
\expandafter\def\csname rep:check:net_worth_mean:public:2025\endcsname{1,232.0}
\expandafter\def\csname rep:check:net_worth_mean:published:2022\endcsname{1,164.1}
\expandafter\def\csname rep:check:net_worth_mean:published:2025\endcsname{1,241.5}
\expandafter\def\csname rep:check:net_worth_median:gap:2022\endcsname{$-$0.10}
\expandafter\def\csname rep:check:net_worth_median:gap:2025\endcsname{+0.09}
\expandafter\def\csname rep:check:net_worth_median:public:2022\endcsname{210.9}
\expandafter\def\csname rep:check:net_worth_median:public:2025\endcsname{216.1}
\expandafter\def\csname rep:check:net_worth_median:published:2022\endcsname{211.1}
\expandafter\def\csname rep:check:net_worth_median:published:2025\endcsname{215.9}
\expandafter\def\csname rep:check:rows:2022\endcsname{22,975}
\expandafter\def\csname rep:check:rows:2025\endcsname{21,800}
\expandafter\def\csname rep:check:weight_total:gap:2022\endcsname{+0.00}
\expandafter\def\csname rep:check:weight_total:gap:2025\endcsname{$-$0.03}
\expandafter\def\csname rep:check:weight_total:public:2022\endcsname{131.3}
\expandafter\def\csname rep:check:weight_total:public:2025\endcsname{135.0}
\expandafter\def\csname rep:check:weight_total:published:2022\endcsname{131.3}
\expandafter\def\csname rep:check:weight_total:published:2025\endcsname{135.0}
\expandafter\def\csname rep:count:all\endcsname{27}
\expandafter\def\csname rep:count:all:replicated\endcsname{16}
\expandafter\def\csname rep:count:primary\endcsname{8}
\expandafter\def\csname rep:count:primary:direction\endcsname{0}
\expandafter\def\csname rep:count:primary:failed\endcsname{0}
\expandafter\def\csname rep:count:primary:inconclusive\endcsname{2}
\expandafter\def\csname rep:count:primary:replicated\endcsname{6}
\expandafter\def\csname rep:count:secondary\endcsname{19}
\expandafter\def\csname rep:count:secondary:direction\endcsname{2}
\expandafter\def\csname rep:count:secondary:failed\endcsname{5}
\expandafter\def\csname rep:count:secondary:inconclusive\endcsname{2}
\expandafter\def\csname rep:count:secondary:replicated\endcsname{10}
\expandafter\def\csname rep:count:status2022:notmet\endcsname{2}
\expandafter\def\csname rep:debt-pinball-candidate:2022\endcsname{49,201}
\expandafter\def\csname rep:debt-pinball-candidate:2022r\endcsname{53,844}
\expandafter\def\csname rep:debt-pinball-candidate:2025\endcsname{53,883}
\expandafter\def\csname rep:debt-pinball-marginal:2022\endcsname{48,277}
\expandafter\def\csname rep:debt-pinball-marginal:2022r\endcsname{52,833}
\expandafter\def\csname rep:debt-pinball-marginal:2025\endcsname{52,793}
\expandafter\def\csname rep:debt-w1-candidate:2022\endcsname{13,001}
\expandafter\def\csname rep:debt-w1-candidate:2022r\endcsname{14,228}
\expandafter\def\csname rep:debt-w1-candidate:2025\endcsname{12,412}
\expandafter\def\csname rep:debt-w1-hotdeck-ratio:2022\endcsname{1.06}
\expandafter\def\csname rep:debt-w1-hotdeck-ratio:2025\endcsname{1.03}
\expandafter\def\csname rep:debt-w1-hotdeck:2022\endcsname{13,770}
\expandafter\def\csname rep:debt-w1-hotdeck:2022r\endcsname{15,070}
\expandafter\def\csname rep:debt-w1-hotdeck:2025\endcsname{12,741}
\expandafter\def\csname rep:debt-w1-ols:2022\endcsname{156,295}
\expandafter\def\csname rep:debt-w1-ols:2022r\endcsname{171,046}
\expandafter\def\csname rep:debt-w1-ols:2025\endcsname{125,966}
\expandafter\def\csname rep:debt-w1-unweighted-ratio:2022\endcsname{2.48}
\expandafter\def\csname rep:debt-w1-unweighted-ratio:2025\endcsname{1.49}
\expandafter\def\csname rep:debt-w1-unweighted:2022\endcsname{32,247}
\expandafter\def\csname rep:debt-w1-unweighted:2022r\endcsname{35,290}
\expandafter\def\csname rep:debt-w1-unweighted:2025\endcsname{18,537}
\expandafter\def\csname rep:debt-zero-candidate:2022\endcsname{0.020}
\expandafter\def\csname rep:debt-zero-candidate:2025\endcsname{0.017}
\expandafter\def\csname rep:debt-zero-forest-max:2022\endcsname{0.020}
\expandafter\def\csname rep:debt-zero-forest-max:2025\endcsname{0.027}
\expandafter\def\csname rep:debt-zero-forest:2022\endcsname{0.018}
\expandafter\def\csname rep:debt-zero-forest:2025\endcsname{0.027}
\expandafter\def\csname rep:debt-zero-quantreg:2022\endcsname{0.151}
\expandafter\def\csname rep:debt-zero-quantreg:2025\endcsname{0.157}
\expandafter\def\csname rep:delta-forest:2022\endcsname{98,496}
\expandafter\def\csname rep:delta-forest:2022r\endcsname{107,792}
\expandafter\def\csname rep:delta-forest:2025\endcsname{29,485}
\expandafter\def\csname rep:delta-unchained:2022\endcsname{22,735}
\expandafter\def\csname rep:delta-unchained:2022r\endcsname{24,881}
\expandafter\def\csname rep:delta-unchained:2025\endcsname{37,883}
\expandafter\def\csname rep:delta-unweighted:2022\endcsname{693,287}
\expandafter\def\csname rep:delta-unweighted:2022r\endcsname{758,718}
\expandafter\def\csname rep:delta-unweighted:2025\endcsname{545,638}
\expandafter\def\csname rep:diag:C:changed\endcsname{1}
\expandafter\def\csname rep:diag:C:changed-ids\endcsname{F2c}
\expandafter\def\csname rep:min-auc-candidate:2022\endcsname{0.732}
\expandafter\def\csname rep:min-auc-candidate:2025\endcsname{0.719}
\expandafter\def\csname rep:min-auc-forest-gap:2022\endcsname{4.3}
\expandafter\def\csname rep:min-auc-forest-gap:2025\endcsname{4.5}
\expandafter\def\csname rep:min-auc-forest-sd:2022\endcsname{+1.83}
\expandafter\def\csname rep:min-auc-forest-sd:2025\endcsname{+2.01}
\expandafter\def\csname rep:min-auc-forest:2022\endcsname{0.775}
\expandafter\def\csname rep:min-auc-forest:2025\endcsname{0.764}
\expandafter\def\csname rep:min-auc-hotdeck-gap:2022\endcsname{2.0}
\expandafter\def\csname rep:min-auc-hotdeck-gap:2025\endcsname{1.1}
\expandafter\def\csname rep:min-auc-hotdeck-sd:2022\endcsname{+0.86}
\expandafter\def\csname rep:min-auc-hotdeck-sd:2025\endcsname{+0.51}
\expandafter\def\csname rep:min-auc-hotdeck:2022\endcsname{0.753}
\expandafter\def\csname rep:min-auc-hotdeck:2025\endcsname{0.731}
\expandafter\def\csname rep:min-auc-marginal:2022\endcsname{0.848}
\expandafter\def\csname rep:min-auc-marginal:2025\endcsname{0.845}
\expandafter\def\csname rep:min-auc-ols:2022\endcsname{0.988}
\expandafter\def\csname rep:min-auc-ols:2025\endcsname{0.983}
\expandafter\def\csname rep:min-auc-quantreg:2022\endcsname{0.927}
\expandafter\def\csname rep:min-auc-quantreg:2025\endcsname{0.907}
\expandafter\def\csname rep:min-auc-unchained-gap:2022\endcsname{0.9}
\expandafter\def\csname rep:min-auc-unchained-gap:2025\endcsname{0.8}
\expandafter\def\csname rep:min-auc-unchained-sd:2022\endcsname{+0.40}
\expandafter\def\csname rep:min-auc-unchained-sd:2025\endcsname{+0.36}
\expandafter\def\csname rep:min-auc-unchained:2022\endcsname{0.742}
\expandafter\def\csname rep:min-auc-unchained:2025\endcsname{0.727}
\expandafter\def\csname rep:min-auc-unweighted-sd:2022\endcsname{$-$0.08}
\expandafter\def\csname rep:min-auc-unweighted-sd:2025\endcsname{$-$0.34}
\expandafter\def\csname rep:min-auc-unweighted:2022\endcsname{0.731}
\expandafter\def\csname rep:min-auc-unweighted:2025\endcsname{0.712}
\expandafter\def\csname rep:min-coverage-candidate:2022\endcsname{0.960}
\expandafter\def\csname rep:min-coverage-candidate:2025\endcsname{0.964}
\expandafter\def\csname rep:min-coverage-hotdeck-sd:2022\endcsname{+0.21}
\expandafter\def\csname rep:min-coverage-hotdeck-sd:2025\endcsname{+0.23}
\expandafter\def\csname rep:min-coverage-hotdeck:2022\endcsname{0.962}
\expandafter\def\csname rep:min-coverage-hotdeck:2025\endcsname{0.965}
\expandafter\def\csname rep:min-coverage-unweighted-sd:2022\endcsname{$-$0.16}
\expandafter\def\csname rep:min-coverage-unweighted-sd:2025\endcsname{+0.39}
\expandafter\def\csname rep:min-coverage-unweighted:2022\endcsname{0.958}
\expandafter\def\csname rep:min-coverage-unweighted:2025\endcsname{0.966}
\expandafter\def\csname rep:min-energy-candidate:2022\endcsname{0.139}
\expandafter\def\csname rep:min-energy-candidate:2025\endcsname{0.098}
\expandafter\def\csname rep:min-energy-hotdeck-sd:2022\endcsname{+0.42}
\expandafter\def\csname rep:min-energy-hotdeck-sd:2025\endcsname{+1.00}
\expandafter\def\csname rep:min-energy-hotdeck:2022\endcsname{0.141}
\expandafter\def\csname rep:min-energy-hotdeck:2025\endcsname{0.102}
\expandafter\def\csname rep:min-energy-unweighted-sd:2022\endcsname{+0.27}
\expandafter\def\csname rep:min-energy-unweighted-sd:2025\endcsname{+0.36}
\expandafter\def\csname rep:min-energy-unweighted:2022\endcsname{0.140}
\expandafter\def\csname rep:min-energy-unweighted:2025\endcsname{0.099}
\expandafter\def\csname rep:min-floor-auc:2022\endcsname{0.496}
\expandafter\def\csname rep:min-floor-auc:2025\endcsname{0.506}
\expandafter\def\csname rep:min-floor-coverage:2022\endcsname{0.987}
\expandafter\def\csname rep:min-floor-coverage:2025\endcsname{0.991}
\expandafter\def\csname rep:min-floor-energy:2022\endcsname{0.009}
\expandafter\def\csname rep:min-floor-energy:2025\endcsname{0.010}
\expandafter\def\csname rep:min-floor-q99:2022\endcsname{1.10}
\expandafter\def\csname rep:min-floor-q99:2025\endcsname{0.990}
\expandafter\def\csname rep:min-q90-candidate:2022\endcsname{1.08}
\expandafter\def\csname rep:min-q90-candidate:2025\endcsname{0.843}
\expandafter\def\csname rep:min-q90-floor:2022\endcsname{1.05}
\expandafter\def\csname rep:min-q90-floor:2025\endcsname{0.991}
\expandafter\def\csname rep:min-q90-hotdeck:2022\endcsname{1.07}
\expandafter\def\csname rep:min-q90-hotdeck:2025\endcsname{0.688}
\expandafter\def\csname rep:min-q90-marginal:2022\endcsname{1.04}
\expandafter\def\csname rep:min-q90-marginal:2025\endcsname{0.999}
\expandafter\def\csname rep:min-q90-unweighted:2022\endcsname{1.43}
\expandafter\def\csname rep:min-q90-unweighted:2025\endcsname{1.09}
\expandafter\def\csname rep:min-q99-candidate:2022\endcsname{0.696}
\expandafter\def\csname rep:min-q99-candidate:2025\endcsname{0.642}
\expandafter\def\csname rep:min-q99-forest:2022\endcsname{0.805}
\expandafter\def\csname rep:min-q99-forest:2025\endcsname{0.699}
\expandafter\def\csname rep:min-q99-hotdeck:2022\endcsname{0.762}
\expandafter\def\csname rep:min-q99-hotdeck:2025\endcsname{1.10}
\expandafter\def\csname rep:min-q99-marginal:2022\endcsname{1.09}
\expandafter\def\csname rep:min-q99-marginal:2025\endcsname{0.994}
\expandafter\def\csname rep:min-q99-unchained:2022\endcsname{0.765}
\expandafter\def\csname rep:min-q99-unchained:2025\endcsname{0.728}
\expandafter\def\csname rep:min-q99-unweighted-sd:2022\endcsname{+6.11}
\expandafter\def\csname rep:min-q99-unweighted-sd:2025\endcsname{+7.52}
\expandafter\def\csname rep:min-q99-unweighted:2022\endcsname{1.93}
\expandafter\def\csname rep:min-q99-unweighted:2025\endcsname{1.79}
\expandafter\def\csname rep:min-tail-candidate:2022\endcsname{0.031}
\expandafter\def\csname rep:min-tail-candidate:2025\endcsname{0.040}
\expandafter\def\csname rep:min-tail-hotdeck-sd:2022\endcsname{+6.24}
\expandafter\def\csname rep:min-tail-hotdeck-sd:2025\endcsname{+3.46}
\expandafter\def\csname rep:min-tail-hotdeck:2022\endcsname{0.051}
\expandafter\def\csname rep:min-tail-hotdeck:2025\endcsname{0.058}
\expandafter\def\csname rep:min-tail-unchained-sd:2022\endcsname{$-$0.63}
\expandafter\def\csname rep:min-tail-unchained-sd:2025\endcsname{+1.13}
\expandafter\def\csname rep:min-tail-unchained:2022\endcsname{0.029}
\expandafter\def\csname rep:min-tail-unchained:2025\endcsname{0.046}
\expandafter\def\csname rep:min-tail-unweighted-ratio:2022\endcsname{5.08}
\expandafter\def\csname rep:min-tail-unweighted-ratio:2025\endcsname{3.12}
\expandafter\def\csname rep:min-tail-unweighted:2022\endcsname{0.159}
\expandafter\def\csname rep:min-tail-unweighted:2025\endcsname{0.125}
\expandafter\def\csname rep:nw-pinball-candidate:2022\endcsname{551,317}
\expandafter\def\csname rep:nw-pinball-candidate:2022r\endcsname{603,350}
\expandafter\def\csname rep:nw-pinball-candidate:2025\endcsname{691,571}
\expandafter\def\csname rep:nw-pinball-marginal:2022\endcsname{443,554}
\expandafter\def\csname rep:nw-pinball-marginal:2022r\endcsname{485,416}
\expandafter\def\csname rep:nw-pinball-marginal:2025\endcsname{531,502}
\expandafter\def\csname rep:nw-w1-candidate:2022\endcsname{135,871}
\expandafter\def\csname rep:nw-w1-candidate:2022r\endcsname{148,694}
\expandafter\def\csname rep:nw-w1-candidate:2025\endcsname{165,931}
\expandafter\def\csname rep:nw-w1-forest-ratio:2022\endcsname{1.72}
\expandafter\def\csname rep:nw-w1-forest-ratio:2025\endcsname{1.18}
\expandafter\def\csname rep:nw-w1-forest:2022\endcsname{234,367}
\expandafter\def\csname rep:nw-w1-forest:2022r\endcsname{256,486}
\expandafter\def\csname rep:nw-w1-forest:2025\endcsname{195,416}
\expandafter\def\csname rep:nw-w1-hotdeck-ratio:2022\endcsname{3.75}
\expandafter\def\csname rep:nw-w1-hotdeck-ratio:2025\endcsname{2.62}
\expandafter\def\csname rep:nw-w1-hotdeck:2022\endcsname{509,013}
\expandafter\def\csname rep:nw-w1-hotdeck:2022r\endcsname{557,053}
\expandafter\def\csname rep:nw-w1-hotdeck:2025\endcsname{434,090}
\expandafter\def\csname rep:nw-w1-ols:2022\endcsname{3,565,923}
\expandafter\def\csname rep:nw-w1-ols:2022r\endcsname{3,902,469}
\expandafter\def\csname rep:nw-w1-ols:2025\endcsname{2,703,593}
\expandafter\def\csname rep:nw-w1-unchained-ratio:2022\endcsname{1.17}
\expandafter\def\csname rep:nw-w1-unchained-ratio:2025\endcsname{1.23}
\expandafter\def\csname rep:nw-w1-unchained:2022\endcsname{158,606}
\expandafter\def\csname rep:nw-w1-unchained:2022r\endcsname{173,575}
\expandafter\def\csname rep:nw-w1-unchained:2025\endcsname{203,814}
\expandafter\def\csname rep:nw-w1-unweighted-ratio:2022\endcsname{6.10}
\expandafter\def\csname rep:nw-w1-unweighted-ratio:2025\endcsname{4.29}
\expandafter\def\csname rep:nw-w1-unweighted:2022\endcsname{829,158}
\expandafter\def\csname rep:nw-w1-unweighted:2022r\endcsname{907,413}
\expandafter\def\csname rep:nw-w1-unweighted:2025\endcsname{711,569}
\expandafter\def\csname rep:pinball:debt:candidate:coded:2022\endcsname{49,201}
\expandafter\def\csname rep:pinball:debt:candidate:coded:2025\endcsname{53,883}
\expandafter\def\csname rep:pinball:debt:candidate:seeds-coded:2022\endcsname{10}
\expandafter\def\csname rep:pinball:debt:candidate:seeds-coded:2025\endcsname{9}
\expandafter\def\csname rep:pinball:debt:candidate:seeds-weighted:2022\endcsname{5}
\expandafter\def\csname rep:pinball:debt:candidate:seeds-weighted:2025\endcsname{7}
\expandafter\def\csname rep:pinball:debt:candidate:weighted:2022\endcsname{48,277}
\expandafter\def\csname rep:pinball:debt:candidate:weighted:2025\endcsname{52,845}
\expandafter\def\csname rep:pinball:debt:hotdeck:coded:2022\endcsname{49,537}
\expandafter\def\csname rep:pinball:debt:hotdeck:coded:2025\endcsname{53,598}
\expandafter\def\csname rep:pinball:debt:hotdeck:seeds-coded:2022\endcsname{10}
\expandafter\def\csname rep:pinball:debt:hotdeck:seeds-coded:2025\endcsname{10}
\expandafter\def\csname rep:pinball:debt:hotdeck:seeds-weighted:2022\endcsname{7}
\expandafter\def\csname rep:pinball:debt:hotdeck:seeds-weighted:2025\endcsname{8}
\expandafter\def\csname rep:pinball:debt:hotdeck:weighted:2022\endcsname{48,279}
\expandafter\def\csname rep:pinball:debt:hotdeck:weighted:2025\endcsname{52,879}
\expandafter\def\csname rep:pinball:debt:marginal:coded:2022\endcsname{48,277}
\expandafter\def\csname rep:pinball:debt:marginal:coded:2025\endcsname{52,793}
\expandafter\def\csname rep:pinball:debt:marginal:seeds-coded:2022\endcsname{0}
\expandafter\def\csname rep:pinball:debt:marginal:seeds-coded:2025\endcsname{0}
\expandafter\def\csname rep:pinball:debt:marginal:seeds-weighted:2022\endcsname{0}
\expandafter\def\csname rep:pinball:debt:marginal:seeds-weighted:2025\endcsname{0}
\expandafter\def\csname rep:pinball:debt:marginal:weighted:2022\endcsname{48,273}
\expandafter\def\csname rep:pinball:debt:marginal:weighted:2025\endcsname{52,832}
\expandafter\def\csname rep:pinball:debt:unweighted:coded:2022\endcsname{50,218}
\expandafter\def\csname rep:pinball:debt:unweighted:coded:2025\endcsname{54,061}
\expandafter\def\csname rep:pinball:debt:unweighted:seeds-coded:2022\endcsname{10}
\expandafter\def\csname rep:pinball:debt:unweighted:seeds-coded:2025\endcsname{10}
\expandafter\def\csname rep:pinball:debt:unweighted:seeds-weighted:2022\endcsname{7}
\expandafter\def\csname rep:pinball:debt:unweighted:seeds-weighted:2025\endcsname{6}
\expandafter\def\csname rep:pinball:debt:unweighted:weighted:2022\endcsname{48,363}
\expandafter\def\csname rep:pinball:debt:unweighted:weighted:2025\endcsname{52,853}
\expandafter\def\csname rep:pinball:networth:candidate:coded:2022\endcsname{551,317}
\expandafter\def\csname rep:pinball:networth:candidate:coded:2025\endcsname{691,571}
\expandafter\def\csname rep:pinball:networth:candidate:seeds-coded:2022\endcsname{10}
\expandafter\def\csname rep:pinball:networth:candidate:seeds-coded:2025\endcsname{10}
\expandafter\def\csname rep:pinball:networth:candidate:seeds-weighted:2022\endcsname{4}
\expandafter\def\csname rep:pinball:networth:candidate:seeds-weighted:2025\endcsname{4}
\expandafter\def\csname rep:pinball:networth:candidate:weighted:2022\endcsname{443,741}
\expandafter\def\csname rep:pinball:networth:candidate:weighted:2025\endcsname{531,304}
\expandafter\def\csname rep:pinball:networth:hotdeck:coded:2022\endcsname{644,921}
\expandafter\def\csname rep:pinball:networth:hotdeck:coded:2025\endcsname{786,929}
\expandafter\def\csname rep:pinball:networth:hotdeck:seeds-coded:2022\endcsname{10}
\expandafter\def\csname rep:pinball:networth:hotdeck:seeds-coded:2025\endcsname{10}
\expandafter\def\csname rep:pinball:networth:hotdeck:seeds-weighted:2022\endcsname{7}
\expandafter\def\csname rep:pinball:networth:hotdeck:seeds-weighted:2025\endcsname{7}
\expandafter\def\csname rep:pinball:networth:hotdeck:weighted:2022\endcsname{443,947}
\expandafter\def\csname rep:pinball:networth:hotdeck:weighted:2025\endcsname{532,123}
\expandafter\def\csname rep:pinball:networth:marginal:coded:2022\endcsname{443,554}
\expandafter\def\csname rep:pinball:networth:marginal:coded:2025\endcsname{531,502}
\expandafter\def\csname rep:pinball:networth:marginal:seeds-coded:2022\endcsname{0}
\expandafter\def\csname rep:pinball:networth:marginal:seeds-coded:2025\endcsname{0}
\expandafter\def\csname rep:pinball:networth:marginal:seeds-weighted:2022\endcsname{0}
\expandafter\def\csname rep:pinball:networth:marginal:seeds-weighted:2025\endcsname{0}
\expandafter\def\csname rep:pinball:networth:marginal:weighted:2022\endcsname{443,736}
\expandafter\def\csname rep:pinball:networth:marginal:weighted:2025\endcsname{531,697}
\expandafter\def\csname rep:pinball:networth:unweighted:coded:2022\endcsname{703,457}
\expandafter\def\csname rep:pinball:networth:unweighted:coded:2025\endcsname{853,940}
\expandafter\def\csname rep:pinball:networth:unweighted:seeds-coded:2022\endcsname{10}
\expandafter\def\csname rep:pinball:networth:unweighted:seeds-coded:2025\endcsname{10}
\expandafter\def\csname rep:pinball:networth:unweighted:seeds-weighted:2022\endcsname{10}
\expandafter\def\csname rep:pinball:networth:unweighted:seeds-weighted:2025\endcsname{7}
\expandafter\def\csname rep:pinball:networth:unweighted:weighted:2022\endcsname{445,725}
\expandafter\def\csname rep:pinball:networth:unweighted:weighted:2025\endcsname{532,547}
\expandafter\def\csname rep:rank:columns\endcsname{15}
\expandafter\def\csname rep:rank:failed\endcsname{1}
\expandafter\def\csname rep:rank:min\endcsname{0.25}
\expandafter\def\csname rep:rank:min:c2st_auc\endcsname{1.00}
\expandafter\def\csname rep:rank:min:energy_distance\endcsname{0.81}
\expandafter\def\csname rep:rank:min:prdc_coverage\endcsname{0.88}
\expandafter\def\csname rep:rank:min:q99_ratio.networth\endcsname{0.67}
\expandafter\def\csname rep:rank:min:w1_over_sd.networth\endcsname{0.98}
\expandafter\def\csname rep:rank:partly\endcsname{4}
\expandafter\def\csname rep:rank:replicated\endcsname{10}
\expandafter\def\csname rep:rank:samebest\endcsname{10}
\expandafter\def\csname rep:rank:scale:c2st_auc\endcsname{0.93}
\expandafter\def\csname rep:rank:scale:energy_distance\endcsname{0.76}
\expandafter\def\csname rep:rank:scale:prdc_coverage\endcsname{0.98}
\expandafter\def\csname rep:rank:scale:q99_ratio.networth\endcsname{0.62}
\expandafter\def\csname rep:rank:scale:w1_over_sd.networth\endcsname{0.57}
\expandafter\def\csname rep:rank:wealth:debt.pinball_loss\endcsname{0.90}
\expandafter\def\csname rep:rank:wealth:debt.wasserstein1\endcsname{0.98}
\expandafter\def\csname rep:rank:wealth:debt.zero_share_error\endcsname{0.25}
\expandafter\def\csname rep:rank:wealth:networth.pinball_loss\endcsname{0.93}
\expandafter\def\csname rep:rank:wealth:networth.wasserstein1\endcsname{0.93}
\expandafter\def\csname rep:repro:cells\endcsname{130}
\expandafter\def\csname rep:repro:identical\endcsname{130}
\expandafter\def\csname rep:repro:maxgap\endcsname{0.00}
\expandafter\def\csname rep:reuse:min:2022:inflated\endcsname{0}
\expandafter\def\csname rep:reuse:min:2022:other-max\endcsname{0.08}
\expandafter\def\csname rep:reuse:min:2022:top-share-max\endcsname{3.9}
\expandafter\def\csname rep:reuse:min:2025:inflated\endcsname{1}
\expandafter\def\csname rep:reuse:min:2025:other-max\endcsname{0.65}
\expandafter\def\csname rep:reuse:min:2025:share-max\endcsname{2.1}
\expandafter\def\csname rep:reuse:min:2025:share-min\endcsname{2.1}
\expandafter\def\csname rep:reuse:min:2025:top-share-max\endcsname{2.5}
\expandafter\def\csname rep:reuse:scale:2022:inflated\endcsname{0}
\expandafter\def\csname rep:reuse:scale:2022:other-max\endcsname{0.09}
\expandafter\def\csname rep:reuse:scale:2022:top-share-max\endcsname{3.4}
\expandafter\def\csname rep:reuse:scale:2025:inflated\endcsname{8}
\expandafter\def\csname rep:reuse:scale:2025:other-max\endcsname{0.04}
\expandafter\def\csname rep:reuse:scale:2025:share-max\endcsname{2.7}
\expandafter\def\csname rep:reuse:scale:2025:share-min\endcsname{1.4}
\expandafter\def\csname rep:reuse:scale:2025:top-share-max\endcsname{3.7}
\expandafter\def\csname rep:scale-auc-candidate:2022\endcsname{0.761}
\expandafter\def\csname rep:scale-auc-candidate:2025\endcsname{0.757}
\expandafter\def\csname rep:scale-auc-forest-gap:2022\endcsname{5.6}
\expandafter\def\csname rep:scale-auc-forest-gap:2025\endcsname{5.9}
\expandafter\def\csname rep:scale-auc-forest-sd:2022\endcsname{+2.63}
\expandafter\def\csname rep:scale-auc-forest-sd:2025\endcsname{+3.03}
\expandafter\def\csname rep:scale-auc-forest:2022\endcsname{0.817}
\expandafter\def\csname rep:scale-auc-forest:2025\endcsname{0.816}
\expandafter\def\csname rep:scale-auc-hotdeck-gap:2022\endcsname{0.6}
\expandafter\def\csname rep:scale-auc-hotdeck-gap:2025\endcsname{$-$0.3}
\expandafter\def\csname rep:scale-auc-hotdeck-sd:2022\endcsname{+0.30}
\expandafter\def\csname rep:scale-auc-hotdeck-sd:2025\endcsname{$-$0.16}
\expandafter\def\csname rep:scale-auc-hotdeck:2022\endcsname{0.768}
\expandafter\def\csname rep:scale-auc-hotdeck:2025\endcsname{0.754}
\expandafter\def\csname rep:scale-auc-marginal:2022\endcsname{0.968}
\expandafter\def\csname rep:scale-auc-marginal:2025\endcsname{0.968}
\expandafter\def\csname rep:scale-auc-ols:2022\endcsname{0.995}
\expandafter\def\csname rep:scale-auc-ols:2025\endcsname{0.993}
\expandafter\def\csname rep:scale-auc-quantreg:2022\endcsname{0.970}
\expandafter\def\csname rep:scale-auc-quantreg:2025\endcsname{0.965}
\expandafter\def\csname rep:scale-auc-unchained-gap:2022\endcsname{5.8}
\expandafter\def\csname rep:scale-auc-unchained-gap:2025\endcsname{4.9}
\expandafter\def\csname rep:scale-auc-unchained-sd:2022\endcsname{+2.70}
\expandafter\def\csname rep:scale-auc-unchained-sd:2025\endcsname{+2.54}
\expandafter\def\csname rep:scale-auc-unchained:2022\endcsname{0.819}
\expandafter\def\csname rep:scale-auc-unchained:2025\endcsname{0.806}
\expandafter\def\csname rep:scale-auc-unweighted-sd:2022\endcsname{$-$0.35}
\expandafter\def\csname rep:scale-auc-unweighted-sd:2025\endcsname{$-$1.12}
\expandafter\def\csname rep:scale-auc-unweighted:2022\endcsname{0.754}
\expandafter\def\csname rep:scale-auc-unweighted:2025\endcsname{0.735}
\expandafter\def\csname rep:scale-coverage-candidate:2022\endcsname{0.969}
\expandafter\def\csname rep:scale-coverage-candidate:2025\endcsname{0.976}
\expandafter\def\csname rep:scale-coverage-hotdeck-sd:2022\endcsname{+0.71}
\expandafter\def\csname rep:scale-coverage-hotdeck-sd:2025\endcsname{$-$0.04}
\expandafter\def\csname rep:scale-coverage-hotdeck:2022\endcsname{0.971}
\expandafter\def\csname rep:scale-coverage-hotdeck:2025\endcsname{0.975}
\expandafter\def\csname rep:scale-coverage-unweighted-sd:2022\endcsname{+0.72}
\expandafter\def\csname rep:scale-coverage-unweighted-sd:2025\endcsname{+0.13}
\expandafter\def\csname rep:scale-coverage-unweighted:2022\endcsname{0.971}
\expandafter\def\csname rep:scale-coverage-unweighted:2025\endcsname{0.976}
\expandafter\def\csname rep:scale-energy-candidate:2022\endcsname{0.118}
\expandafter\def\csname rep:scale-energy-candidate:2025\endcsname{0.085}
\expandafter\def\csname rep:scale-energy-hotdeck-sd:2022\endcsname{+0.24}
\expandafter\def\csname rep:scale-energy-hotdeck-sd:2025\endcsname{+1.19}
\expandafter\def\csname rep:scale-energy-hotdeck:2022\endcsname{0.119}
\expandafter\def\csname rep:scale-energy-hotdeck:2025\endcsname{0.089}
\expandafter\def\csname rep:scale-energy-unweighted-sd:2022\endcsname{+0.57}
\expandafter\def\csname rep:scale-energy-unweighted-sd:2025\endcsname{+0.20}
\expandafter\def\csname rep:scale-energy-unweighted:2022\endcsname{0.120}
\expandafter\def\csname rep:scale-energy-unweighted:2025\endcsname{0.086}
\expandafter\def\csname rep:scale-floor-auc:2022\endcsname{0.499}
\expandafter\def\csname rep:scale-floor-auc:2025\endcsname{0.510}
\expandafter\def\csname rep:scale-floor-coverage:2022\endcsname{0.986}
\expandafter\def\csname rep:scale-floor-coverage:2025\endcsname{0.990}
\expandafter\def\csname rep:scale-floor-energy:2022\endcsname{0.012}
\expandafter\def\csname rep:scale-floor-energy:2025\endcsname{0.013}
\expandafter\def\csname rep:scale-floor-q99:2022\endcsname{1.10}
\expandafter\def\csname rep:scale-floor-q99:2025\endcsname{0.990}
\expandafter\def\csname rep:scale-q90-candidate:2022\endcsname{1.000}
\expandafter\def\csname rep:scale-q90-candidate:2025\endcsname{0.751}
\expandafter\def\csname rep:scale-q90-floor:2022\endcsname{1.05}
\expandafter\def\csname rep:scale-q90-floor:2025\endcsname{0.991}
\expandafter\def\csname rep:scale-q90-hotdeck:2022\endcsname{1.07}
\expandafter\def\csname rep:scale-q90-hotdeck:2025\endcsname{0.761}
\expandafter\def\csname rep:scale-q90-marginal:2022\endcsname{1.06}
\expandafter\def\csname rep:scale-q90-marginal:2025\endcsname{0.988}
\expandafter\def\csname rep:scale-q90-unweighted:2022\endcsname{1.34}
\expandafter\def\csname rep:scale-q90-unweighted:2025\endcsname{1.04}
\expandafter\def\csname rep:scale-q99-candidate:2022\endcsname{0.767}
\expandafter\def\csname rep:scale-q99-candidate:2025\endcsname{0.630}
\expandafter\def\csname rep:scale-q99-debt-candidate:2022\endcsname{1.08}
\expandafter\def\csname rep:scale-q99-debt-candidate:2025\endcsname{0.957}
\expandafter\def\csname rep:scale-q99-debt-unweighted:2022\endcsname{1.57}
\expandafter\def\csname rep:scale-q99-debt-unweighted:2025\endcsname{1.32}
\expandafter\def\csname rep:scale-q99-fin-candidate:2022\endcsname{0.650}
\expandafter\def\csname rep:scale-q99-fin-candidate:2025\endcsname{0.644}
\expandafter\def\csname rep:scale-q99-fin-unweighted:2022\endcsname{1.29}
\expandafter\def\csname rep:scale-q99-fin-unweighted:2025\endcsname{1.15}
\expandafter\def\csname rep:scale-q99-forest:2022\endcsname{1.15}
\expandafter\def\csname rep:scale-q99-forest:2025\endcsname{0.916}
\expandafter\def\csname rep:scale-q99-hotdeck:2022\endcsname{0.747}
\expandafter\def\csname rep:scale-q99-hotdeck:2025\endcsname{2.85}
\expandafter\def\csname rep:scale-q99-marginal:2022\endcsname{1.12}
\expandafter\def\csname rep:scale-q99-marginal:2025\endcsname{1.00}
\expandafter\def\csname rep:scale-q99-nfin-candidate:2022\endcsname{0.772}
\expandafter\def\csname rep:scale-q99-nfin-candidate:2025\endcsname{0.811}
\expandafter\def\csname rep:scale-q99-nfin-unweighted:2022\endcsname{3.08}
\expandafter\def\csname rep:scale-q99-nfin-unweighted:2025\endcsname{1.71}
\expandafter\def\csname rep:scale-q99-unchained:2022\endcsname{0.699}
\expandafter\def\csname rep:scale-q99-unchained:2025\endcsname{0.692}
\expandafter\def\csname rep:scale-q99-unweighted-sd:2022\endcsname{+7.32}
\expandafter\def\csname rep:scale-q99-unweighted-sd:2025\endcsname{+7.50}
\expandafter\def\csname rep:scale-q99-unweighted:2022\endcsname{2.25}
\expandafter\def\csname rep:scale-q99-unweighted:2025\endcsname{1.78}
\expandafter\def\csname rep:scale-tail-candidate:2022\endcsname{0.027}
\expandafter\def\csname rep:scale-tail-candidate:2025\endcsname{0.048}
\expandafter\def\csname rep:scale-tail-hotdeck-sd:2022\endcsname{$-$0.90}
\expandafter\def\csname rep:scale-tail-hotdeck-sd:2025\endcsname{+10.50}
\expandafter\def\csname rep:scale-tail-hotdeck:2022\endcsname{0.024}
\expandafter\def\csname rep:scale-tail-hotdeck:2025\endcsname{0.103}
\expandafter\def\csname rep:scale-tail-unchained-sd:2022\endcsname{+0.50}
\expandafter\def\csname rep:scale-tail-unchained-sd:2025\endcsname{+2.74}
\expandafter\def\csname rep:scale-tail-unchained:2022\endcsname{0.028}
\expandafter\def\csname rep:scale-tail-unchained:2025\endcsname{0.062}
\expandafter\def\csname rep:scale-tail-unweighted-ratio:2022\endcsname{4.94}
\expandafter\def\csname rep:scale-tail-unweighted-ratio:2025\endcsname{1.69}
\expandafter\def\csname rep:scale-tail-unweighted:2022\endcsname{0.133}
\expandafter\def\csname rep:scale-tail-unweighted:2025\endcsname{0.081}
\expandafter\def\csname rep:seeds-debt-w1-unweighted:2022\endcsname{10}
\expandafter\def\csname rep:seeds-debt-w1-unweighted:2025\endcsname{9}
\expandafter\def\csname rep:seeds-debt-zero-forest:2022\endcsname{4}
\expandafter\def\csname rep:seeds-debt-zero-forest:2025\endcsname{8}
\expandafter\def\csname rep:seeds-nw-w1-forest:2022\endcsname{6}
\expandafter\def\csname rep:seeds-nw-w1-forest:2025\endcsname{6}
\expandafter\def\csname rep:seeds-nw-w1-hotdeck:2022\endcsname{9}
\expandafter\def\csname rep:seeds-nw-w1-hotdeck:2025\endcsname{9}
\expandafter\def\csname rep:seeds-nw-w1-unchained:2022\endcsname{7}
\expandafter\def\csname rep:seeds-nw-w1-unchained:2025\endcsname{8}
\expandafter\def\csname rep:seeds-nw-w1-unweighted:2022\endcsname{10}
\expandafter\def\csname rep:seeds-nw-w1-unweighted:2025\endcsname{10}
\expandafter\def\csname rep:seeds-scale-auc-forest:2022\endcsname{10}
\expandafter\def\csname rep:seeds-scale-auc-forest:2025\endcsname{9}
\expandafter\def\csname rep:seeds-scale-auc-unchained:2022\endcsname{9}
\expandafter\def\csname rep:seeds-scale-auc-unchained:2025\endcsname{9}
\expandafter\def\csname rep:units:candidate\endcsname{4.8}
\expandafter\def\csname rep:units:forest-max\endcsname{7.3}
\expandafter\def\csname rep:units:quantreg\endcsname{7.7}


Every \scf{} experiment in the body of the paper uses the 2022 wave, and
every design decision in the estimator and the harness was made with 2022
results in view. The Federal Reserve released the 2025 wave in October 2026
\citep{fed2026scf}. We reran all three \scf{} experiments on it --- the
within-\scf{} task and both profiles of the population view --- with the same
code, locked environment, seeds, protocol, and receiver files (ASEC 2025 for
the minimal profile, pooled ASEC 2023--2025 at populace scale); only the donor
wave changed. Before any method was fit on the new wave we committed a file
fixing which of the paper's findings would count as replicated and by what
rule. This appendix reports the result under those rules as written.

\paragraph{The 2025 file.} The public summary extract holds
\rep{check:families:2025} families in five implicates. Every predictor and
target the two profiles use is present with the same codes; the one changed
definition is race, where the Fed now folds the Middle Eastern or North
African response into the white non-Hispanic class \citep{fed2026scf}.
(The separate Hispanic-origin questions dropped from the 2025 questionnaire
never entered this variable.) The weights sum to
\rep{check:weight_total:public:2025} million families, against
\rep{check:weight_total:published:2025} million in the Fed's report, and the
identity networth = fin + nfin $-$ debt holds on every row. Pooled over
implicates, all-family median and mean net worth are
\rep{check:net_worth_median:public:2025} and
\rep{check:net_worth_mean:public:2025} thousand dollars, against published
values of \rep{check:net_worth_median:published:2025} and
\rep{check:net_worth_mean:published:2025} thousand
(gaps of \rep{check:net_worth_median:gap:2025} and
\rep{check:net_worth_mean:gap:2025} percent). The Fed's tables use internal
data; the public file is disclosure-limited and omits the families whose net
worth reaches the Forbes 400 threshold, seven in each wave
\citep{fed2026scfcodebook}. The 2022 public
file's gaps are \rep{check:net_worth_median:gap:2022} and
\rep{check:net_worth_mean:gap:2022} percent.

\paragraph{Dollars.} The 2025 extract is in 2025 dollars. On the same day the
Fed replaced the 2022 extract at its published address with a copy re-based
to 2025 dollars (every dollar column multiplied by the CPI-U-RS ratio
4789/4376; weights and codes unchanged). The paper's 2022 runs used the
earlier, 2022-dollar file, and the repository's loader now accepts only files
it can identify by hash, so the 2022 results regenerate regardless of which
copy is downloaded. The comparisons below are scale-free or within-wave
ratios. Dollar levels in the verdict tables are in each wave's own dollars.

\paragraph{Rules.} Within the \scf{}, a claim that method B is worse than
method A is met when B is worse on average and on at least eight of ten
paired seeds. On the population view the noise unit is the sampling floor's
seed-to-seed standard deviation: a difference is met beyond two floor
standard deviations (with eight seeds agreeing), and a tie is met within one.
Where the paper states an effect size, it holds if the 2025 effect is within
a factor of two of the 2022 effect. Each rule is also applied to the 2022
runs; \rep{count:status2022:notmet} of the paper's 2022 claims do not meet
these rules on the 2022 data themselves, and for those the best possible 2025 verdict is inconclusive.

\paragraph{Verdicts.}
Tables~\ref{tab:scf2025-verdicts-wealth}--\ref{tab:scf2025-verdicts-view-b}
list all \rep{count:all} findings. Of the
\rep{count:primary} that the abstract or conclusion states:
\rep{count:primary:replicated} replicated, \rep{count:primary:inconclusive}
inconclusive, \rep{count:primary:failed} not replicated. Of the
\rep{count:secondary} others: \rep{count:secondary:replicated} replicated,
\rep{count:secondary:direction} replicated in direction with an effect smaller
than the band allows, \rep{count:secondary:inconclusive} inconclusive,
\rep{count:secondary:failed} not replicated.

\paragraph{What replicated.} The weighting result holds. Fit unweighted, the
estimator is \rep{nw-w1-unweighted-ratio:2025} times worse on net-worth $W_1$
within the \scf{} (\rep{nw-w1-unweighted-ratio:2022} in 2022), on
\rep{seeds-nw-w1-unweighted:2025} of ten seeds, and on the population view its
imputed net-worth q99 is \rep{min-q99-unweighted:2025} and
\rep{scale-q99-unweighted:2025} times the holdout's
(\rep{min-q99-unweighted:2022} and \rep{scale-q99-unweighted:2022} in 2022).
Weighting is again the largest of the three ablation effects. Chaining again
binds only at scale. With two targets the unchained ablation's AUC sits
\rep{min-auc-unchained-sd:2025} floor standard deviations from the
candidate's; on the four-component balance sheet unchaining costs
\rep{scale-auc-unchained-gap:2025} AUC points
(\rep{scale-auc-unchained-gap:2022} in 2022), on
\rep{seeds-scale-auc-unchained:2025} of ten seeds, and the ungated forest
loses \rep{scale-auc-forest-gap:2025} points. The classifier separates the
unconditional marginal draw as sharply as before (AUC
\rep{scale-auc-marginal:2025} against the candidate's
\rep{scale-auc-candidate:2025} at scale), and within the \scf{} the same draw
again posts the best pinball loss on both targets. That last result is a
property of the metric: the pinball loss reads the imputed quantiles
unweighted and scores them against weighted truth. Read under receiver
weights, the candidate's and the marginal draw's net-worth losses in 2025 are
\rep{pinball:networth:candidate:weighted:2025} and
\rep{pinball:networth:marginal:weighted:2025} (as coded,
\rep{pinball:networth:candidate:coded:2025} and
\rep{pinball:networth:marginal:coded:2025}), and the marginal draw is the
lower on \rep{pinball:networth:candidate:seeds-weighted:2025} of ten seeds,
not \rep{pinball:networth:candidate:seeds-coded:2025}; the 2022 wave shows
the same (\rep{pinball:networth:candidate:seeds-weighted:2022} of ten, not
\rep{pinball:networth:candidate:seeds-coded:2022}). The body's pinball
columns are left as published. Hot deck stays within one
floor standard deviation of the candidate on AUC in both profiles
(\rep{min-auc-hotdeck-sd:2025} and \rep{scale-auc-hotdeck-sd:2025}). The
candidate still understates the extreme tail: its net-worth q99 ratio is
\rep{min-q99-candidate:2025} and \rep{scale-q99-candidate:2025}.

\paragraph{What did not.} \emph{Hot deck's tail at scale} (F4c, F5b). In 2022
hot deck matched the candidate on the net-worth tail at scale: a q99 ratio of
\rep{scale-q99-hotdeck:2022} against the candidate's
\rep{scale-q99-candidate:2022}, and a $W_1$/sd of
\rep{scale-tail-hotdeck:2022} against \rep{scale-tail-candidate:2022}. In
2025 its q99 ratio at scale is \rep{scale-q99-hotdeck:2025} and its $W_1$/sd
\rep{scale-tail-hotdeck:2025}, against the candidate's
\rep{scale-q99-candidate:2025} and \rep{scale-tail-candidate:2025}. Donor
reuse produces the q99 inflation. On
\rep{reuse:scale:2025:inflated} of the ten seeds, one donor above the
holdout's q99 is matched to between \rep{reuse:scale:2025:share-min} and
\rep{reuse:scale:2025:share-max} percent of receiver weight. A donor
whose receivers carry more than one percent of the weight, by more than any
single receiver's share, puts the imputed q99 at or above its own net worth,
so that match alone lifts the imputed q99 above the holdout's. On the remaining seeds the largest such share is
\rep{reuse:scale:2025:other-max} percent, and in 2022 it never exceeded
\rep{reuse:scale:2022:other-max} percent. In the minimal profile hot deck's
mean q99 ratio is \rep{min-q99-hotdeck:2025}, which looks like a correct
tail and is not one: \rep{reuse:min:2025:inflated} of the ten seeds is
inflated the same way, by a donor carrying \rep{reuse:min:2025:share-max}
percent of receiver weight. The body's statement that hot deck slightly beats the candidate on
the tail at scale holds for 2022 only.

\emph{The candidate as closest to the floor} (F3c). At scale hot deck's mean AUC
is \rep{scale-auc-hotdeck:2025}, below the candidate's
\rep{scale-auc-candidate:2025}, so the candidate is not the lowest among
itself and the baselines. The gap is \rep{scale-auc-hotdeck-sd:2025} floor
standard deviations: on this wave the candidate and hot deck are tied for
closest to the floor.

\emph{Unchaining and the tail} (F2c). At scale the unchained ablation's
net-worth $W_1$/sd is \rep{scale-tail-unchained:2025} against the candidate's
\rep{scale-tail-candidate:2025}, \rep{scale-tail-unchained-sd:2025} floor
standard deviations apart; \rep{F2c:parts-met:2025} of the
\rep{F2c:parts} ties behind ``essentially unchanged'' were met.

\emph{Gating and debt within the \scf{}} (G1). The 2022 runs showed no gating
advantage on the debt zero share. In 2025 the ungated forest's error is
\rep{debt-zero-forest:2025} against the candidate's
\rep{debt-zero-candidate:2025}, and it is the worse of the two on
\rep{seeds-debt-zero-forest:2025} of ten seeds. The 2022 tie did not recur;
the difference that replaced it favours the gated estimator.

\emph{Smaller effects.} Unweighted is \rep{debt-w1-unweighted-ratio:2025}
times worse on debt $W_1$ within the \scf{}, against
\rep{debt-w1-unweighted-ratio:2022} in 2022. On the population view its tail
$W_1$/sd is \rep{min-tail-unweighted-ratio:2025} and
\rep{scale-tail-unweighted-ratio:2025} times the candidate's, against
\rep{min-tail-unweighted-ratio:2022} and \rep{scale-tail-unweighted-ratio:2022};
the body's ``five times'' does not hold at scale. Part of that is the
candidate's own tail: its $W_1$/sd at scale rose from
\rep{scale-tail-candidate:2022} to \rep{scale-tail-candidate:2025}.

\emph{Inconclusive.} Each inconclusive primary finding is a claim that
several ties hold at once, and in each, one tie fell between one and two
floor standard deviations. For the unweighted ablation (F1a), \rep{F1a:parts-met:2025} of
\rep{F1a:parts} ties were met: at scale its AUC is
\rep{scale-auc-unweighted:2025} against the candidate's
\rep{scale-auc-candidate:2025} (\rep{scale-auc-unweighted-sd:2025} floor
standard deviations; the unweighted fit is the harder to tell from the
holdout). For hot deck (F4a), \rep{F4a:parts-met:2025} of \rep{F4a:parts}
were met: at scale its energy distance is \rep{scale-energy-hotdeck:2025}
against \rep{scale-energy-candidate:2025}
(\rep{scale-energy-hotdeck-sd:2025}). The claim that the ungated forest is
worse than the candidate on net-worth $W_1$ within the \scf{} stayed within
seed noise, as it was in 2022 (\rep{seeds-nw-w1-forest:2025} of ten seeds in
2025, \rep{seeds-nw-w1-forest:2022} in 2022). The claim that removing
chaining raises net-worth $W_1$ within the \scf{} met its rule in 2025 (a
factor of \rep{nw-w1-unchained-ratio:2025}, on
\rep{seeds-nw-w1-unchained:2025} of ten seeds) but not in 2022
(\rep{seeds-nw-w1-unchained:2022} of ten), so under the cap stated above it
stays inconclusive.

\paragraph{Outside the rules.} Energy distance fell for every method (the
candidate's from \rep{min-energy-candidate:2022} to
\rep{min-energy-candidate:2025} in the minimal profile and from
\rep{scale-energy-candidate:2022} to \rep{scale-energy-candidate:2025} at
scale). The criteria anticipated a shift shared by all methods, because
\scf{} 2025 incomes refer to the same year as the ASEC 2025 receiver's; we
did not test that explanation. The paper makes no claim about the net-worth
q90 ratio. On this wave the candidate understates it at scale:
\rep{scale-q90-candidate:2025}, against \rep{scale-q90-candidate:2022} in
2022 and a floor of \rep{scale-q90-floor:2025}.

\paragraph{Method rankings.} Ranking the eight displayed methods in each
column of the three \scf{} tables, \rep{rank:replicated} of
\rep{rank:columns} rankings replicated under the pre-registered
rank-correlation rule, \rep{rank:partly} partly, and \rep{rank:failed} not
(debt zero-share error within the \scf{}, where the candidate, its
ablations, the ungated forest, hot deck, and the marginal draw sit within a
seed standard deviation of one another in both waves). The best method is unchanged in
\rep{rank:samebest} columns.

\paragraph{Diagnostics.} Rerunning the 2022 experiments in the environment
used for 2025 reproduced the committed results exactly:
\rep{repro:identical} of \rep{repro:cells} displayed table cells match on all
ten seeds. Code, environment, and seeds are therefore held fixed between the
two columns of the tables below; the donor wave, and with it the dollar
basis, is what differs. Rerunning the 2022 wave in 2025 dollars left every
2022 status unchanged except \rep{diag:C:changed-ids}. (On the population
view that rerun also shifts the donor's income predictors against the
receiver's, whose incomes are not re-based.) The same arm shows that the forest-based
methods are not exactly invariant to the unit of account: a pure change of
units within the \scf{} leaves OLS and the marginal draw unchanged and moves
the forest-based methods' ten-seed $W_1$ means by up to
\rep{units:forest-max} percent. One new wave is one more draw. Where 2022 and
2025 disagree, these runs cannot say which is the unusual year.

\begin{table}[p]
\centering
\scriptsize
% Generated by `imp replication` -- do not edit by hand. Pre-registered rules of docs/scf-2025-replication-criteria.md applied to the committed 2022 and 2025 runs. An asterisk marks a primary finding.
\begin{tabular}{@{}lp{0.20\textwidth}p{0.26\textwidth}p{0.26\textwidth}p{0.12\textwidth}@{}}
\hline
ID & Finding & 2022 & 2025 & Verdict \\
\hline
W1$^{\ast}$ & Fit unweighted, the same estimator is six times worse on net-worth W1. & 829,158 vs 135,871; ratio 6.10; worse on 10/10 seeds & 711,569 vs 165,931; ratio 4.29; worse on 10/10 seeds & Replicated \\
W2 & Unweighted is 2.5 times worse on debt W1. & 32,247 vs 13,001; ratio 2.48; worse on 10/10 seeds & 18,537 vs 12,412; ratio 1.49; worse on 9/10 seeds & Replicated in direction, magnitude differs (smaller) \\
W3a & Weighting is the largest of the three ablation effects on net-worth W1. & deltas: unweighted +693,287; ungated forest +98,496; unchained +22,735 & deltas: unweighted +545,638; ungated forest +29,485; unchained +37,883 & Replicated \\
W3b & The ungated, unchained forest is worse than the candidate on net-worth W1 (1.7 times). & 234,367 vs 135,871; ratio 1.72; worse on 6/10 seeds & 195,416 vs 165,931; ratio 1.18; worse on 6/10 seeds & Inconclusive \\
W3c & Removing chaining raises net-worth W1 by 17 percent. & 158,606 vs 135,871; ratio 1.17; worse on 7/10 seeds & 203,814 vs 165,931; ratio 1.23; worse on 8/10 seeds & Inconclusive \\
W4a & Hot deck nearly ties the candidate on debt W1. & 13,770 vs 13,001; ratio 1.06; hot deck worse on 5/10 seeds, better on 5/10 & 12,741 vs 12,412; ratio 1.03; hot deck worse on 5/10 seeds, better on 5/10 & Replicated \\
W4b & Hot deck is 3.7 times worse on net-worth W1. & 509,013 vs 135,871; ratio 3.75; worse on 9/10 seeds & 434,090 vs 165,931; ratio 2.62; worse on 9/10 seeds & Replicated \\
W5 & OLS is catastrophic on both targets. & OLS W1 is 12.0x (debt) and 26.2x (net worth) the candidate's; highest on both & OLS W1 is 10.1x (debt) and 16.3x (net worth) the candidate's; highest on both & Replicated \\
W6 & Linear quantile regression leaves 15 percent of the debt zero mass unaccounted for, against at most 2 percent for forest-based methods. & quantile regression 0.151; forest-based 0.012 to 0.020 & quantile regression 0.157; forest-based 0.017 to 0.027 & Replicated \\
W7 & The unconditional weighted marginal draw has the best pinball loss on both targets. & debt: marginal 48,277, candidate 49,201, lowest: marginal draw; networth: marginal 443,554, candidate 551,317, lowest: marginal draw & debt: marginal 52,793, candidate 53,883, lowest: marginal draw; networth: marginal 531,502, candidate 691,571, lowest: marginal draw & Replicated \\
G1 & Gating shows no advantage on SCF debt zero share. & ungated forest 0.018 vs candidate 0.020; forest worse on 4/10 seeds, better on 6/10 & ungated forest 0.027 vs candidate 0.017; forest worse on 8/10 seeds, better on 2/10 & Not replicated \\
\hline
\end{tabular}

\caption{Replication verdicts, within-\scf{} wealth task. Each cell gives the
numbers the rule was applied to. Dollar amounts are in each wave's own
dollars; ratios and seed counts are comparable across waves. An asterisk
marks a finding the abstract or conclusion states.}
\label{tab:scf2025-verdicts-wealth}
\end{table}

\begin{table}[p]
\centering
\scriptsize
% Generated by `imp replication` -- do not edit by hand. Pre-registered rules of docs/scf-2025-replication-criteria.md applied to the committed 2022 and 2025 runs. An asterisk marks a primary finding.
\begin{tabular}{@{}lp{0.20\textwidth}p{0.26\textwidth}p{0.26\textwidth}p{0.12\textwidth}@{}}
\hline
ID & Finding & 2022 & 2025 & Verdict \\
\hline
F0 & The sampling floor behaves as designed. & minimal: AUC 0.496, energy 0.009, coverage 0.987, q99 ratio 1.10; scale: AUC 0.499, energy 0.012, coverage 0.986, q99 ratio 1.10 & minimal: AUC 0.506, energy 0.010, coverage 0.991, q99 ratio 0.99; scale: AUC 0.510, energy 0.013, coverage 0.990, q99 ratio 0.99 & Replicated \\
F1a$^{\ast}$ & The unweighted ablation ties the candidate on energy distance, coverage, and C2ST AUC in both profiles. & largest gap 0.72 floor sd (coverage, scale: 0.971 vs 0.969); 6/6 ties met & largest gap 1.12 floor sd (AUC, scale: 0.735 vs 0.757); 5/6 ties met & Inconclusive \\
F1b$^{\ast}$ & Its imputed net-worth q99 is about twice the holdout's (1.93 and 2.25). & minimal: 1.93 vs candidate 0.70 (6.1 floor sd; 10/10 seeds); scale: 2.25 vs candidate 0.77 (7.3 floor sd; 10/10 seeds) & minimal: 1.79 vs candidate 0.64 (7.5 floor sd; 10/10 seeds); scale: 1.78 vs candidate 0.63 (7.5 floor sd; 10/10 seeds) & Replicated \\
F1c & Its tail W1/sd is five times the candidate's. & minimal: 0.159 vs 0.031, ratio 5.08 (10/10 seeds); scale: 0.133 vs 0.027, ratio 4.94 (10/10 seeds) & minimal: 0.125 vs 0.040, ratio 3.12 (9/10 seeds); scale: 0.081 vs 0.048, ratio 1.69 (8/10 seeds) & Replicated in direction, magnitude differs (smaller) \\
F2a$^{\ast}$ & With two targets the unchained ablation is indistinguishable from the candidate. & 0.742 vs 0.732; gap 0.40 floor sd & 0.727 vs 0.719; gap 0.36 floor sd & Replicated \\
F2b$^{\ast}$ & On the four-component balance sheet, fitting targets independently costs 5.8 AUC points. & 0.819 vs 0.761; gap +0.058 (2.70 floor sd); higher on 9/10 seeds & 0.806 vs 0.757; gap +0.049 (2.54 floor sd); higher on 9/10 seeds & Replicated \\
F2c & Unchaining leaves the marginal and tail statistics essentially unchanged at scale. & largest gap 0.76 floor sd (coverage, scale: 0.967 vs 0.969); 4/4 ties met & largest gap 2.74 floor sd (W1/sd, scale: 0.062 vs 0.048); 3/4 ties met & Not replicated \\
F2d & The ungated forest loses 5.6 AUC points at scale. & 0.817 vs 0.761; gap +0.056 (2.63 floor sd); higher on 10/10 seeds & 0.816 vs 0.757; gap +0.059 (3.03 floor sd); higher on 9/10 seeds & Replicated \\
\hline
\end{tabular}

\caption{Replication verdicts, population view: the floor, weighting, and
chaining. ``Floor sd'' is the seed-to-seed standard deviation of the
sampling-floor row. An asterisk marks a finding the abstract or conclusion
states.}
\label{tab:scf2025-verdicts-view}
\end{table}

\begin{table}[p]
\centering
\scriptsize
% Generated by `imp replication` -- do not edit by hand. Pre-registered rules of docs/scf-2025-replication-criteria.md applied to the committed 2022 and 2025 runs. An asterisk marks a primary finding.
\begin{tabular}{@{}lp{0.20\textwidth}p{0.26\textwidth}p{0.26\textwidth}p{0.12\textwidth}@{}}
\hline
ID & Finding & 2022 & 2025 & Verdict \\
\hline
F3a$^{\ast}$ & At scale the classifier separates the unconditional marginal draw almost perfectly (AUC 0.97 against the candidate's 0.76). & 0.968 vs 0.761; gap +0.207 (9.68 floor sd); higher on 10/10 seeds & 0.968 vs 0.757; gap +0.211 (10.90 floor sd); higher on 10/10 seeds & Replicated \\
F3b & Conditioning richness makes the test decisive: the marginal draw's AUC rises from the minimal profile to scale. & 0.848 (minimal) to 0.968 (scale): rise +0.120 & 0.845 (minimal) to 0.968 (scale): rise +0.123 & Replicated \\
F3c & The candidate stays closest to the floor. & minimal: candidate 0.732, lowest baseline 0.753; scale: candidate 0.761, lowest baseline 0.768 & minimal: candidate 0.719, lowest baseline 0.731; scale: candidate 0.757, lowest baseline 0.754 & Not replicated \\
F4a$^{\ast}$ & Hot deck is within noise of the candidate on energy distance and coverage in both profiles. & largest gap 0.71 floor sd (coverage, scale: 0.971 vs 0.969); 4/4 ties met & largest gap 1.19 floor sd (energy, scale: 0.089 vs 0.085); 3/4 ties met & Inconclusive \\
F4b$^{\ast}$ & Hot deck is within 2.1 and 0.7 AUC points of the candidate. & largest gap 0.86 floor sd (AUC, minimal: 0.753 vs 0.732); 2/2 ties met & largest gap 0.51 floor sd (AUC, minimal: 0.731 vs 0.719); 2/2 ties met & Replicated \\
F4c & At scale hot deck slightly beats the candidate on tail W1/sd. & hot deck 0.024 vs candidate 0.027 ($-$0.0028, 0.90 floor sd) & hot deck 0.103 vs candidate 0.048 (+0.0550, 10.50 floor sd) & Not replicated \\
F5a & The candidate understates the extreme tail (net-worth q99 ratio 0.70 and 0.77). & minimal: 0.70, below 1 on 10/10 seeds; scale: 0.77, below 1 on 9/10 seeds & minimal: 0.64, below 1 on 10/10 seeds; scale: 0.63, below 1 on 10/10 seeds & Replicated \\
F5b & Hot deck does not fix the tail either, while the marginal draw sits near the floor (read from the 2022 tables, not the text). & minimal: hot deck 0.76 vs candidate 0.70 (0.33 floor sd); marginal 1.09 vs floor 1.10 (0.05 floor sd); scale: hot deck 0.75 vs candidate 0.77 (0.10 floor sd); marginal 1.12 vs floor 1.10 (0.07 floor sd) & minimal: hot deck 1.10 vs candidate 0.64 (3.01 floor sd); marginal 0.99 vs floor 0.99 (0.02 floor sd); scale: hot deck 2.85 vs candidate 0.63 (14.56 floor sd); marginal 1.00 vs floor 0.99 (0.09 floor sd) & Not replicated \\
\hline
\end{tabular}

\caption{Replication verdicts, population view, continued: the marginal
draw, hot deck, and the tail.}
\label{tab:scf2025-verdicts-view-b}
\end{table}

\begin{table}[p]
\centering
\footnotesize
\resizebox{\textwidth}{!}{% Generated by `imp tables` -- do not edit by hand. SCF wealth task, SCF 2025 donor, mean (sd) over 10 paired seeds.
\begin{tabular}{lrrrrr}
\hline
Method & Debt pinball & Debt $W_1$ & Debt zero-share err. & Net worth pinball & Net worth $W_1$ \\
\hline
populace-fit & 53,883 (2,516) & \textbf{12,412} (4,467) & \textbf{0.017} (0.022) & 691,571 (65,700) & \textbf{165,931} (90,856) \\
\quad -- unweighted & 54,061 (2,224) & 18,537 (7,279) & 0.021 (0.018) & 853,940 (69,647) & 711,569 (313,876) \\
\quad -- unchained & 53,883 (2,516) & \textbf{12,412} (4,467) & \textbf{0.017} (0.022) & 693,429 (46,890) & 203,814 (57,212) \\
\quad -- ungated/unchained forest & 55,184 (1,927) & 13,575 (3,675) & 0.027 (0.011) & 689,111 (21,625) & 195,416 (62,104) \\
OLS & 69,765 (6,754) & 125,966 (31,799) & 0.230 (0.013) & 1,116,887 (197,719) & 2,703,593 (964,714) \\
Quantile regression & 54,264 (1,963) & 21,064 (5,644) & 0.157 (0.015) & 672,364 (41,739) & 235,061 (46,789) \\
NND hot deck (py-statmatch) & 53,598 (2,172) & 12,741 (3,353) & 0.020 (0.020) & 786,929 (52,769) & 434,090 (251,013) \\
Weighted marginal draw & \textbf{52,793} (2,339) & 13,178 (3,605) & 0.019 (0.010) & \textbf{531,502} (32,912) & 215,488 (90,663) \\
\hline
\end{tabular}
}
\caption{Table~\ref{tab:scf-wealth} on the 2025 wave (2025 dollars).}
\label{tab:scf-wealth-2025}
\end{table}

\begin{table}[p]
\centering
\footnotesize
\resizebox{\textwidth}{!}{% Generated by `imp tables` -- do not edit by hand. SCF->CPS population view vs held-out SCF, SCF 2025 donor, mean (sd) over 10 donor splits. AUC of 0.5 is indistinguishable; q99 ratio of 1 is a perfect tail. Minimal profile: 6 shared predictors, 2 targets, ASEC 2025 receiver.
\begin{tabular}{lrrrrr}
\hline
Method & Energy distance & Coverage & C2ST AUC & NW $W_1$/sd & NW q99 ratio \\
\hline
SCF sample (floor) & 0.010 (0.004) & 0.991 (0.005) & 0.506 (0.022) & 0.018 (0.005) & 0.990 (0.153) \\
\hline
populace-fit & \textbf{0.098} (0.019) & 0.964 (0.009) & 0.719 (0.009) & 0.040 (0.010) & 0.642 (0.122) \\
\quad -- unweighted & 0.099 (0.019) & \textbf{0.966} (0.009) & \textbf{0.712} (0.010) & 0.125 (0.068) & 1.79 (0.726) \\
\quad -- unchained & 0.098 (0.019) & 0.964 (0.010) & 0.727 (0.012) & 0.046 (0.017) & 0.728 (0.170) \\
\quad -- ungated/unchained forest & 0.100 (0.018) & 0.961 (0.013) & 0.764 (0.013) & 0.039 (0.008) & 0.699 (0.176) \\
OLS & 0.116 (0.025) & 0.759 (0.067) & 0.983 (0.004) & 0.351 (0.130) & 0.676 (0.215) \\
Quantile regression & 0.101 (0.020) & 0.957 (0.012) & 0.907 (0.007) & 0.060 (0.010) & 0.462 (0.091) \\
NND hot deck (py-statmatch) & 0.102 (0.020) & 0.965 (0.009) & 0.731 (0.011) & 0.058 (0.036) & 1.10 (0.808) \\
Weighted marginal draw & 0.101 (0.019) & 0.940 (0.017) & 0.845 (0.010) & \textbf{0.021} (0.008) & \textbf{0.994} (0.181) \\
\hline
\end{tabular}
}
\caption{Table~\ref{tab:harness} on the 2025 wave: the population view,
minimal profile, with the \scf{} 2025 donor.}
\label{tab:harness-2025}
\end{table}

\begin{table}[p]
\centering
\footnotesize
\resizebox{\textwidth}{!}{% Generated by `imp tables` -- do not edit by hand. SCF->CPS population view vs held-out SCF, SCF 2025 donor, mean (sd) over 10 donor splits. AUC of 0.5 is indistinguishable; q99 ratio of 1 is a perfect tail. Populace-scale profile: 10 shared predictors, 4 chained targets, pooled ASEC 2023-2025 receiver.
\begin{tabular}{lrrrrr}
\hline
Method & Energy distance & Coverage & C2ST AUC & NW $W_1$/sd & NW q99 ratio \\
\hline
SCF sample (floor) & 0.013 (0.003) & 0.990 (0.006) & 0.510 (0.019) & 0.018 (0.005) & 0.990 (0.153) \\
\hline
populace-fit & 0.085 (0.015) & 0.976 (0.005) & 0.757 (0.048) & 0.048 (0.016) & 0.630 (0.165) \\
\quad -- unweighted & 0.086 (0.015) & \textbf{0.976} (0.004) & \textbf{0.735} (0.013) & 0.081 (0.042) & 1.78 (0.745) \\
\quad -- unchained & \textbf{0.085} (0.015) & 0.973 (0.006) & 0.806 (0.014) & 0.062 (0.059) & 0.692 (0.195) \\
\quad -- ungated/unchained forest & 0.092 (0.017) & 0.968 (0.009) & 0.816 (0.011) & 0.030 (0.007) & 0.916 (0.261) \\
OLS & 0.103 (0.018) & 0.902 (0.025) & 0.993 (0.002) & 0.349 (0.130) & 0.663 (0.203) \\
Quantile regression & 0.087 (0.015) & 0.960 (0.011) & 0.965 (0.004) & 0.064 (0.011) & 0.425 (0.074) \\
NND hot deck (py-statmatch) & 0.089 (0.016) & 0.975 (0.005) & 0.754 (0.009) & 0.103 (0.043) & 2.85 (1.30) \\
Weighted marginal draw & 0.089 (0.015) & 0.930 (0.009) & 0.968 (0.003) & \textbf{0.019} (0.007) & \textbf{1.00} (0.147) \\
\hline
\end{tabular}
}
\caption{Table~\ref{tab:harness-scale} on the 2025 wave: the population view,
populace-scale profile, with the \scf{} 2025 donor.}
\label{tab:harness-scale-2025}
\end{table}

\clearpage
